\chapter{List of references}

1. B. Harkness; The Seedlist Handbook, compiled and updated by M. G. Harkness
and D. D'Angelo,Timber Press, Portland, Oregon, 1986

2. J. D. Bewley and M. Black; Seeds, Physiology of Development and Germination,
Plenum Press, New York, 1985

3. A. Khan, editor; Physiology and Biochemistry of Seed Dormancy and
Germination, Elsevier Publishing Co., New York, 1977

4. A. M. Mayer and A. Poljakoff-Mayber; The Germination of Seeds, 3rd edition,
Pergamon Press, New York, 1982; ref. 4A, 4th edition, 1988

5. K. E. Ovcharov; Physiological Basis of Seed Germination, trans. from the
Russian, Amerind Publishing Co., New Delhi, India, original 1969, trans. 1977

6. D. A. Priestley; Seed Aging, Constock Publishing Associates, Cornell Univeristy
Press, Ithaca and London, 1986

7. M. G. Nikolaeva; Physiology of Deep Dormancy in Seeds, trans. from Russian,
Israel Program Scientific Translations, Jerusalem, 1969

8. V. S. Summerhayes, Wild Orchids of Britain, Collins, London, 1951

9. R. Stewart and B. Presley, Bulletin of the American Rock Garden Society, 49, 253,
1991

10. M. G. Nikolaeva and Coworkers, Bot. Zh. Leningrad fl, 61(1987)

11. J. W. Bradbeer, Seed Dormancy and Germination, Blackie, New York, 1988

12. Ref. 2, pp. 201 and 334; Ref. 5, pp. 46 and 334; Ref. 6, P. 33

13. Ref. 2, p. 202

14. Ref. 7, p. 16

15. Ref. 3, p. 32; Ref. 5, p. 168

16. Ref. 4, 4th edition, 1988, pp. 54-55

17. W. Kinzel, Frost und Licht, Neue Tabellen, Eugen Ulmer, Stuttgart, 1926

18. H. A. Borkwith and Coworkers, Proc. NatI. Acad. Sd. 3fl, 662 (1952)

19. Phytochrome, K. Mitrakos and W. Shropshire, Jr., Editors, Academic Press, NY,
1972

20. Recent Advances in the Develo pment and Germination of Seeds, P. B. Taylorson,
Plenum Press, New York, 1989, Ch. by J. D. Smith and Coworkers, p. 57

21. J. A. DeGreef and Coworkers, p. 241 in Ref. 20

22. Ref. 7, pp. 41-45 .

23. Gibberelins, N. Takahashi, B. 0. Phinney, and J. MacMillan, Eds., Springer-
Verlag, New York, 1991

24. Abscisic Acid, F. T. Addicott, Ed., Praeger Publishers, New York, 1983

25. P. Stokes, Encyclopedia of Plant Physiology 15, 746-803 (1965)

26. David Lynn, Ann. Rev. Plant Phys. Al, 497-526 (1990)

27. M. A. Cohn in ref. 20, p. 61

28. L. V. Barton, Physiology of Seeds, Chronica Botanica Co., Waltham, Mass.,
1953,

29. Ref. 4, pp. 158-161

30. A. M. J. Else and D. N. Riemer, J. Aquat. Plant Manage. 22, 22-25, 1984

31. L. V. Barton and W. Crocker, Twenty Years of Seed Research, Faber and Faber,
London, no publication date, Chapter 2

32. Ref. 3, p. 30; Ref. 4, p. 50; and Ref. 5, p.75

33. Ref. 4, p.21; Ref. 6, pp. 77 and 84 -

34. Ref. 2, p. 93; Ref. 6, p. 51

35. Ref. 2, p. 97; Ref. 6, p. 59

36. Ref. 2, pp. 93 and 342; Ref. 6, p. 21

37. N. Deno, Surface Sterilization Method, Bull. Am. Rock Garden Soc., spring
1986, p. 1

38. Stephen Jay Gould, Discover, pp. 40-44, December 1992

39. P. L. Privalov, Advances in Protein Chemistry 3, 172 (1979); C. Tanford,
Advances in Protein Chemistry, 121 (1968); I. W. Sizer. Enzyme
Kinetics 3, 35(1943);

40. M. Dixon and E. C. Webb,Enzvmes, Longman Group Ltd., London, 3rd edition,
Chapter IV, pp. 164-182

41. R. Jaynes, J. Am. Hort. Sci., 9, 608, (1971)

42. N. C. Deno, The Mystery of Lilium Tenuifolium, Bulletin of the American Rock
Garden Society, a 225, 1992

43. N. C. Deno, Unraveling the Secrets of Primula Kisoana, Bulletin of the American
Rock Garden Society, 5Q, 211, 1992

\chapter{Digest of symbols and abbreviations}

\section{Letters}

\begin{itemize}
\item[DS] dry storage with the temperature (40 or 70) given after the DS
WC washed in water and cleaned daily usually for a period of seven days
GA-3 is gibberelic acid-3.
\item[d] days
\item[m] months
\item[y] years
\item[T] temperature in degrees Fahrenheit
\end{itemize}

\section{Time Numbers}

\begin{itemize}
\item[D-70] A species in which the seed requires dry storage before the seeds
will germinate at 70.
\item[40-70] means that the seeds were subjected to 3 m at 40 followed by
germination at 70. This system can be extended indefinitely.
\item[70-40-70-40] means that the seeds were subjected to 3 m at 70, 3 m at
\item[40, 3 m at 70] and shifted to 40 whereupon germination began.
\item[70D] means that the seeds were in the dark.
\item[70L] means that the seeds were in the light as described in Chapter 3.
\item[70 GA-3] means that the seeds were treated with GA-3 as described in Chapter 3.
\item[40\%] In referring to germination this means 40\% of the normal sized seed coats
germinated. Some attempts were made to discard and not count undersize
seed coats, seed coats that readily crushed and were obviously empty, and
overly thin seed coats that were empty of endosperm.
\item[in 6-10w] Referring to germination this means that germination started at the end of
the 6th week and ended at the end of the 10th w.
\end{itemize}

\section{Abbreviations}

\begin{itemize}
\item[ind. t] is the induction time,which is the time between the seeds being shifted or
started in a particular cycle and the time the seeds start to germinate in that
same cycle.
\item[250/old] means that 25\% of the seeds germinate in each day once germination starts.
half life This is the time for half of the seeds to germinate again starting from the time
germination starts
\item[zero order rate] This is a rate in which a constant amount germinates in each time
period
\item[first order rate] This is a rate in which a constant fraction germinates in each time
period
\item[germ.] germinated
\end{itemize}

